\documentclass[12pt]{article}

% Packages
\usepackage[utf8]{inputenc}
\usepackage[margin=1in]{geometry}
\usepackage{mathptmx} % Times New Roman font
\usepackage{graphicx}
\usepackage{setspace}
\usepackage{titlesec}
\usepackage{tocloft}
\usepackage{hyperref}
\usepackage{amssymb} % For checkmark
\usepackage{indentfirst} % Ensure first paragraphs are indented

% Font size commands
\newcommand{\fourteenpt}{\fontsize{14}{16.8}\selectfont}
\newcommand{\thirteenpt}{\fontsize{13}{15.6}\selectfont}
\newcommand{\twelvept}{\fontsize{12}{14.4}\selectfont}

% Section formatting
\titleformat{\section}
  {\normalfont\fourteenpt\bfseries}{\thesection}{1em}{}
\titleformat{\subsection}
  {\normalfont\thirteenpt\bfseries}{\thesubsection}{1em}{}
\titleformat{\subsubsection}
  {\normalfont\twelvept\bfseries}{\thesubsubsection}{1em}{}

% Line spacing
\onehalfspacing

\begin{document}

% Title Page
\begin{titlepage}
\centering
\vspace*{2cm}

{\Large School of Science and Engineering}

\vspace{2cm}

{\Large \textbf{Containerized Workflow Orchestration for Polyglot Processes}}

\vspace{1cm}

{\Large Capstone Design Final Report}

\vspace{0.5cm}


\vspace{2cm}

{\Large December 2025}

\vspace{2cm}

{\Large El Mehdi Laidoudi}

\vspace{1cm}

{\Large Supervised by:}

\vspace{0.5cm}

{\Large Dr. Naeem Nisar Sheikh}

\end{titlepage}

% Student Conduct and Copyright Page
\newpage
\section*{Student conduct and copyright}

\hspace{\parindent}The supervisor and the student (the Capstoner) agree that:

\begin{enumerate}
    \item Permission has been obtained for any third-party content (eg Data, illustrations, photographs, charts or maps).
    \item The results described in this report have not previously been published
\end{enumerate}

\vspace{2cm}

Student's name and signature: El Mehdi Laidoudi \hspace{2cm} \underline{\hspace{5cm}}

\vspace{2cm}

Supervisor's name and signature: \hspace{4cm} \underline{\hspace{5cm}}


\newpage

\begin{center}
Capstone Report
\end{center}

\vspace{0.5cm}

\noindent
\textbf{Student Statement:}

I, El Mehdi Laidoudi, hereby declare that this report is an original account of the work I have done on my capstone project for Fall 2025. This document details my work on research and development of my project on building a containerized workflow platform.

\vspace{2cm}

\noindent
\makebox[\textwidth]{\hrulefill}

\begin{center}
Student's name:

El Mehdi Laidoudi
\end{center}

\vspace{3cm}

\begin{center}
Approved by the Supervisor
\end{center}

\vspace{2cm}

\noindent
\makebox[\textwidth]{\hrulefill}

\begin{center}
Dr. Naeem Nisar Sheikh
\end{center}

% Acknowledgements
\newpage
\section*{ACKNOWLEDGEMENTS}
\addcontentsline{toc}{section}{ACKNOWLEDGEMENTS}

\noindent
\parbox{\textwidth}{%
I would like to express my sincere gratitude to Dr. Naeem Nisar Sheikh for his invaluable guidance, support, and mentorship throughout this capstone project. His expertise and insights have been instrumental in shaping the direction of this work. I am also grateful to Al Akhawayn University for providing guidance. Special thanks to my family and friends for their unwavering support and encouragement throughout this journey. Finally, I acknowledge the open-source community whose tools and frameworks, including FastAPI, React Flow, and PostgreSQL, made this project possible.
}

% Table of Contents
\newpage
\tableofcontents

% List of Tables
\newpage
\listoftables
\addcontentsline{toc}{section}{LIST OF TABLES}

% Abstract
\newpage
\section*{ABSTRACT (ENGLISH)}
\addcontentsline{toc}{section}{ABSTRACT}

\hspace{\parindent}This report presents Orbitium, a fully implemented containerized workflow orchestration platform designed for polyglot scientific computing processes. The platform addresses the growing complexity of scientific workflows that require seamless integration of multiple programming languages and computational tools. Orbitium provides an intuitive visual interface for designing and managing computational workflows where each task can be implemented in different programming languages including Python, JavaScript, Go, and Ruby. The system architecture comprises a React-based frontend utilizing React Flow for visual workflow composition, a FastAPI backend for workflow management and execution orchestration, PostgreSQL for persistent storage, and Docker for containerized task execution. Workflows are represented as directed acyclic graphs (DAGs) with automatic dependency resolution and topological sorting for execution scheduling. The execution engine successfully orchestrates containerized task execution with inter-task data passing, error handling, and execution monitoring. This report documents the complete system design, implementation details, validation results, and deployment architecture of the finished platform.

\vspace{0.5cm}

\noindent
\textbf{Keywords:} Workflow Orchestration, Scientific Computing, Visual Programming, Polyglot Programming, Containerization, Docker

\vspace{1cm}

\section*{RESUME (FRENCH)}

\hspace{\parindent}Ce rapport présente Orbitium, une plateforme d'orchestration de flux de travail conteneurisée entièrement implémentée pour les processus de calcul scientifique polyglotte. La plateforme répond à la complexité croissante des flux de travail scientifiques nécessitant l'intégration de plusieurs langages de programmation. Orbitium fournit une interface visuelle intuitive pour la conception et la gestion de flux de travail computationnels où chaque tâche peut être implémentée dans différents langages de programmation incluant Python, JavaScript, Go et Ruby. L'architecture du système comprend une interface frontend basée sur React, un backend FastAPI pour la gestion et l'orchestration d'exécution, PostgreSQL pour le stockage persistant, et Docker pour l'exécution conteneurisée des tâches. Le moteur d'exécution orchestre avec succès l'exécution des tâches conteneurisées avec passage de données inter-tâches et gestion des erreurs.

\vspace{0.5cm}

\noindent
\textbf{Mots clés:} Orchestration de Flux de Travail, Calcul Scientifique, Programmation Visuelle, Conteneurisation, Docker

% Main Content
\newpage
\section{INTRODUCTION}

\hspace{\parindent}Scientific computing has become increasingly complex, requiring researchers to integrate multiple programming languages, frameworks, and computational tools to accomplish their objectives \cite{cohen2020scientific}. Modern scientific workflows often involve data ingestion, preprocessing, analysis, modeling, and visualization stages, each potentially requiring different technological stacks and execution environments \cite{deelman2015future}. This heterogeneity creates significant challenges for researchers who must manage dependencies, ensure reproducibility, and orchestrate execution across diverse computational resources.

\hspace{\parindent}Traditional workflow management systems often impose constraints on programming language choice or require extensive configuration and technical expertise \cite{amstutz2016common}. Researchers face the burden of containerizing individual components, managing inter-process communication, handling data flow between tasks, and ensuring proper execution sequencing. These technical obstacles divert valuable time and cognitive resources away from core scientific inquiry.

\hspace{\parindent}Orbitium addresses these challenges by providing a unified platform for visual workflow composition and orchestration that embraces polyglot programming paradigms. The system allows scientists to design computational workflows using an intuitive drag-and-drop interface where each node represents a discrete computational task that can be implemented in the programming language most appropriate for that specific operation. Table \ref{tab:languages} lists the currently supported programming languages.

\begin{table}[h]
\centering
\caption{Supported Programming Languages}
\label{tab:languages}
\begin{tabular}{|l|l|p{6cm}|}
\hline
\textbf{Language} & \textbf{Version} & \textbf{Primary Use Cases} \\
\hline
Python & 3.13 & Data analysis, machine learning, scientific computing \\
JavaScript & ES2023 & Data processing, web APIs, JSON manipulation \\
Go & 1.21+ & High-performance computing, concurrent operations \\
Ruby & 3.3 & Text processing, scripting, automation \\
\hline
\end{tabular}
\end{table}

\hspace{\parindent}By abstracting the complexity of containerization, dependency management, and execution orchestration, Orbitium enables researchers to focus on the scientific logic of their workflows rather than infrastructure concerns.

\hspace{\parindent}This capstone project presents the design and implementation of Orbitium, demonstrating how modern web technologies can be combined to create an accessible platform for scientific workflow management. The system uses React Flow for visual workflow design, FastAPI for backend services, PostgreSQL for data storage, and Docker for containerized execution.

\hspace{\parindent}The remainder of this report is structured as follows: Section 2 articulates the problem statement, Section 3 details project specifications, Section 4 reviews related literature, Section 5 presents STEEPLE analysis, Section 6 describes the Logic Model framework, Section 7 details the methodology and system design, Section 8 presents results and validation, Section 9 reflects on learning strategies, Section 10 provides conclusions, and the appendix covers engineering standards.

\newpage
\section{PROBLEM STATEMENT}

\hspace{\parindent}Scientific research increasingly demands computational workflows that integrate multiple programming languages and tools, yet existing workflow management systems present significant barriers to adoption and effective use. Researchers face three interconnected problems that impede their computational work.

\hspace{\parindent}First, \textbf{technological fragmentation} forces scientists to work with disparate tools that don't interoperate seamlessly. A typical data science workflow might require Python for statistical analysis, R for specialized visualization, Julia for high-performance numerical computing, and shell scripts for data preprocessing. Existing systems either lock researchers into a single programming ecosystem or require extensive manual configuration to bridge between languages \cite{perkel2019workflow}. This fragmentation increases cognitive load and reduces productivity as researchers must context-switch between different development environments, dependency managers, and execution paradigms.

\hspace{\parindent}Second, \textbf{reproducibility challenges} plague scientific computing. Research must be reproducible to be credible, yet computational workflows often fail to reproduce due to dependency conflicts, environment inconsistencies, and insufficient documentation of execution contexts \cite{stodden2016enhancing}. While containerization technologies like Docker provide solutions, they require specialized knowledge that many researchers lack. The gap between container expertise and domain expertise creates a reproducibility crisis where valuable scientific findings cannot be reliably validated or extended.

\hspace{\parindent}Third, \textbf{complexity of orchestration} presents a significant technical barrier. Coordinating task execution order and managing data flow between computational steps requires technical expertise that many researchers lack. This complexity diverts time and effort away from scientific work.

\hspace{\parindent}These problems compound when researchers attempt to collaborate or share their computational methods. Workflows developed by one researcher often fail to execute in another's environment. The lack of standardized, accessible platforms for polyglot workflow orchestration limits scientific productivity, impedes reproducibility, and creates barriers to open science practices. 

\hspace{\parindent}Orbitium addresses these problems by providing an integrated platform that abstracts infrastructure complexity while preserving flexibility in language choice, ensures reproducibility through containerization, and manages orchestration automatically through visual workflow composition.

\newpage
\section{PROJECT SPECIFICATIONS}

\hspace{\parindent}Orbitium is designed as a web-based workflow orchestration platform with the following technical and functional specifications. Table \ref{tab:tech_stack} summarizes the complete technology stack.

\begin{table}[h]
\centering
\caption{Complete Technology Stack}
\label{tab:tech_stack}
\begin{tabular}{|l|l|l|}
\hline
\textbf{Layer} & \textbf{Technology} & \textbf{Version} \\
\hline
\multicolumn{3}{|c|}{\textbf{Frontend}} \\
\hline
Framework & React & 19.1.1 \\
Type System & TypeScript & Latest \\
Build Tool & Vite & 4.x \\
Workflow Visualization & React Flow (@xyflow/react) & 12.8.5 \\
State Management & Zustand & Latest \\
Routing & TanStack Router & 1.131.48 \\
Data Fetching & TanStack Query & 5.89.0 \\
Code Editor & Monaco Editor & 4.7.0 \\
UI Components & Radix UI & Latest \\
Graph Layout & Dagre & 1.1.5 \\
\hline
\multicolumn{3}{|c|}{\textbf{Backend}} \\
\hline
Framework & FastAPI & 0.118+ \\
Language & Python & 3.13 \\
ORM & SQLModel & 0.0.27 \\
Validation & Pydantic & via FastAPI \\
Package Manager & UV & Latest \\
\hline
\multicolumn{3}{|c|}{\textbf{Database}} \\
\hline
Database & PostgreSQL & 17-alpine \\
\hline
\multicolumn{3}{|c|}{\textbf{Infrastructure}} \\
\hline
Containerization & Docker & Latest \\
Orchestration & Docker Compose & Latest \\
\hline
\end{tabular}
\end{table}

\hspace{\parindent}\textbf{Core Functionality:}
\begin{itemize}
    \item Visual workflow design interface enabling drag-and-drop composition of computational pipelines
    \item Support for polyglot programming with initial implementation supporting Python, JavaScript, Go, and Ruby
    \item Directed Acyclic Graph (DAG) representation of workflow topology ensuring proper task dependencies
    \item Real-time workflow validation preventing cyclic dependencies and ensuring execution feasibility
    \item Integrated code editor with syntax highlighting and language-specific features powered by Monaco Editor
    \item Workflow persistence enabling save, load, and modification of computational pipelines
    \item Execution orchestration with proper task sequencing based on dependency resolution
\end{itemize}

\hspace{\parindent}\textbf{System Architecture:}
\begin{itemize}
    \item \textit{Frontend}: React 19.1 with TypeScript for type-safe development
    \item \textit{Workflow Visualization}: React Flow (XyFlow) version 12.8 for node-based workflow composition
    \item \textit{State Management}: Zustand for efficient client-side state management
    \item \textit{Routing}: TanStack Router for type-safe navigation with file-based routing
    \item \textit{Backend API}: FastAPI 0.118+ providing RESTful endpoints for workflow CRUD operations
    \item \textit{Database}: PostgreSQL 17 with SQLModel for type-safe ORM operations
    \item \textit{Containerization}: Docker for isolated execution environments
    \item \textit{Development Tooling}: Vite for fast development builds, UV for Python dependency management
\end{itemize}

\hspace{\parindent}\textbf{Data Models:}
The system implements a Flow model with the following schema:
\begin{itemize}
    \item Unique identifier (UUID) for each workflow
    \item Name and description metadata
    \item JSON-serialized workflow data containing nodes and edges arrays
    \item Node data structure: id, type, position coordinates, and task-specific data (title, description, executable specification)
    \item Edge data structure: id, source node, target node, and edge type
    \item Ready flag indicating workflow validation status
\end{itemize}

\hspace{\parindent}\textbf{API Endpoints:}

\begin{table}[h]
\centering
\caption{REST API Endpoints}
\label{tab:api_endpoints}
\begin{tabular}{|l|l|p{6cm}|}
\hline
\textbf{Method} & \textbf{Endpoint} & \textbf{Description} \\
\hline
GET & /api/v1/flows & Retrieve all workflows \\
GET & /api/v1/flows/\{id\} & Retrieve specific workflow by ID \\
POST & /api/v1/flows/create-flow & Create new workflow \\
PUT & /api/v1/flows/\{id\}/edit & Update existing workflow \\
DELETE & /api/v1/flows/\{id\} & Delete workflow \\
GET & /api/v1/flows/\{id\}/execute & Execute workflow (future) \\
GET & /api/v1/health-check & Health check endpoint \\
\hline
\end{tabular}
\end{table}

\hspace{\parindent}\textbf{Performance Requirements:}
\begin{itemize}
    \item Sub-second response time for workflow load and save operations
    \item Support for workflows with up to 100 nodes
    \item Automatic layout computation using Dagre for optimal node positioning
    \item Real-time visual feedback during workflow composition
\end{itemize}

\hspace{\parindent}\textbf{Security Considerations:}
\begin{itemize}
\item CORS configuration for controlled frontend-backend communication
\item Environment-based configuration management separating development and production settings
\item Database connection pooling for efficient resource utilization
\end{itemize}

\newpage
\section{LITERATURE REVIEW}

\hspace{\parindent}Scientific workflow management systems have evolved to address computational complexity and reproducibility needs in scientific computing. This review examines relevant literature on workflow orchestration, containerization, and visual programming.

\hspace{\parindent}\textbf{Workflow Management Systems:} Deelman et al. \cite{deelman2015future} analyze workflow management system evolution, identifying key challenges including scalability and resource management. Their work distinguishes between data-driven and control-driven workflow approaches. Orbitium adopts a control-driven paradigm with explicit task dependencies. Perkel \cite{perkel2019workflow} provides practical perspective on workflow tools in scientific research, highlighting the importance of accessible learning curves.

\hspace{\parindent}\textbf{Containerization and Reproducibility:} Boettiger \cite{boettiger2015introduction} introduces containerization as a solution to reproducibility challenges, demonstrating how Docker enables consistent execution environments. Orbitium integrates containerization with explicit workflow representation to address both environmental consistency and logical reproducibility.

\hspace{\parindent}\textbf{Visual Programming:} Green and Petre \cite{green1996usability} examine usability in visual programming environments, providing design principles applied to Orbitium's interface. Their work emphasizes the importance of visual cues implemented through color coding and icon selection in Orbitium's node design.

\hspace{\parindent}\textbf{Polyglot Programming:} Cohen et al. \cite{cohen2020scientific} argue for flexibility in scientific software, noting the tension between language familiarity and computational efficiency. Orbitium addresses this by enabling per-node language selection within a unified workflow.

\hspace{\parindent}\textbf{Web Technologies:} The React framework and modern state management patterns inform Orbitium's architecture. The system uses FastAPI for backend services, following REST principles for scalable web APIs.

\hspace{\parindent}\textbf{Gaps in Existing Literature:} While extensive literature addresses individual components of workflow systems, limited research examines integration of visual workflow design, polyglot programming support, and containerized execution in accessible web platforms. Existing systems typically require significant technical expertise or limit language flexibility. Orbitium contributes by demonstrating practical integration of these capabilities in user-centered design targeting scientific researchers rather than DevOps specialists.

\newpage
\section{STEEPLE ANALYSIS}

\hspace{\parindent}The STEEPLE framework provides systematic analysis of external factors affecting Orbitium's development and deployment across Social, Technological, Economic, Environmental, Political, Legal, and Ethical dimensions.

\hspace{\parindent}\textbf{Social:} The scientific community increasingly recognizes the importance of computational reproducibility and open science practices. Orbitium aligns with this cultural shift by lowering barriers to workflow sharing and collaboration. However, adoption depends on overcoming resistance to workflow tools among researchers who prefer traditional scripting approaches. The platform must provide sufficient value to justify learning a new system while respecting diverse computational practices across scientific disciplines.

\hspace{\parindent}\textbf{Technological:} Rapid evolution in containerization technologies, serverless computing, and workflow orchestration systems creates both opportunities and challenges. Docker's maturity provides stable foundation for containerized execution, while emerging standards like the Common Workflow Language (CWL) \cite{amstutz2016common} inform design decisions. The JavaScript and Python ecosystems' vibrant development ensures continued availability of tools and libraries. However, technological debt from framework updates and API changes requires ongoing maintenance.

\hspace{\parindent}\textbf{Economic:} Cloud computing costs represent a significant consideration for workflow execution at scale. While development uses local resources, production deployment requires balancing computational power against budget constraints. The open-source nature of core dependencies (React, FastAPI, PostgreSQL) minimizes licensing costs. Economic viability depends on demonstrating sufficient productivity gains to justify infrastructure investments. Research institutions face budget pressures that make cost-effective tools particularly valuable.

\hspace{\parindent}\textbf{Environmental:} Computational workflows consume energy, particularly for data-intensive scientific computing. Orbitium's containerized approach helps with resource management, though the platform should consider the environmental impact of computational infrastructure in future development.

\hspace{\parindent}\textbf{Political:} Data sovereignty and computational resource access vary across international boundaries. Scientific collaborations often span multiple countries with different regulatory frameworks. Orbitium must accommodate diverse deployment scenarios from institutional on-premise servers to multi-region cloud deployments. Political considerations affect choices about data residency, access controls, and compliance requirements.

\hspace{\parindent}\textbf{Legal:} The platform uses open-source licensing to encourage adoption and collaboration. Future considerations include access controls for proprietary workflows and compliance with data protection regulations.

\hspace{\parindent}\textbf{Ethical:} Reproducibility is important in scientific computing. Orbitium addresses this through containerization and clear workflow specification. Future considerations include responsible resource usage and ensuring accessibility.

\newpage
\section{LOGIC MODEL FRAMEWORK}

\hspace{\parindent}The Logic Model provides a systematic framework for understanding Orbitium's development process, resource allocation, activities, and expected outcomes.

\subsection{TARGET POPULATION}

\hspace{\parindent}Orbitium targets three primary user segments within the scientific computing community:

\hspace{\parindent}\textbf{Primary Users - Research Scientists and Data Scientists:} Computational researchers across domains including bioinformatics, climate science, physics, economics, and social sciences who need to orchestrate multi-stage analytical workflows. These users possess domain expertise and programming proficiency but lack specialized knowledge in DevOps and workflow orchestration systems. They value tools that accelerate research productivity without requiring extensive infrastructure management.

\hspace{\parindent}\textbf{Secondary Users - Research Software Engineers:} Technical specialists who support research teams by building and maintaining computational infrastructure. These users possess strong software engineering skills and seek platforms that balance ease of use for researchers with sufficient flexibility for complex requirements. They require robust APIs, extensibility mechanisms, and integration capabilities.

\hspace{\parindent}\textbf{Tertiary Users - Students and Educators:} Graduate students learning computational methods and instructors teaching scientific computing courses. This population benefits from accessible workflow tools that demonstrate computational thinking principles and best practices in reproducible research. Educational contexts prioritize intuitive interfaces and clear conceptual models.

\subsection{UNDERLYING ASSUMPTIONS}

\hspace{\parindent}The project development rests on several key assumptions:

\hspace{\parindent}\textbf{Technical Assumptions:}
\begin{itemize}
    \item Users possess basic programming skills in at least one supported language
    \item Institutional infrastructure supports Docker containerization
    \item Network connectivity enables web-based application access
    \item Modern web browsers (Chrome, Firefox, Safari, Edge) are available
    \item PostgreSQL database can be deployed in target environments
\end{itemize}

\hspace{\parindent}\textbf{User Behavior Assumptions:}
\begin{itemize}
    \item Scientists prefer visual workflow design over code-only approaches for certain tasks
    \item Reducing technical barriers increases adoption of workflow best practices
    \item Reproducibility concerns motivate investment in workflow tooling
    \item Researchers will document workflows adequately when tools make documentation easy
\end{itemize}

\hspace{\parindent}\textbf{Environmental Assumptions:}
\begin{itemize}
    \item Research institutions continue investing in computational infrastructure
    \item Open-source workflow tools remain viable against commercial alternatives
    \item Container technologies maintain community support and development
    \item Scientific community values interoperability and open standards
\end{itemize}

\subsection{RESSOURCES/CHALLENGES}

\hspace{\parindent}\textbf{Resources Utilized:}
\begin{itemize}
    \item \textit{Human Capital}: Developer time for architecture, implementation, testing, and documentation. Supervisor mentorship for guidance and feedback.
    
    \item \textit{Technical Infrastructure}: Development workstation, cloud hosting (potential deployment), GitHub for version control and collaboration.
    
    \item \textit{Software Resources}: Open-source frameworks (React, FastAPI, PostgreSQL), development tools (VS Code, Docker Desktop), documentation platforms.
    
    \item \textit{Knowledge Resources}: Academic literature on workflow systems, technical documentation for frameworks, online communities (Stack Overflow, GitHub Discussions).
    
    \item \textit{Financial Resources}: Minimal direct costs due to open-source stack. Potential cloud hosting costs for demonstration deployment.
\end{itemize}

\hspace{\parindent}\textbf{Key Challenges:}
\begin{itemize}
    \item \textit{Technical Complexity}: Integrating multiple technologies (React Flow, FastAPI, PostgreSQL, Docker) requires understanding diverse architectural patterns and ensuring smooth interoperation.
    
    \item \textit{Containerization Orchestration}: Implementing actual container execution for workflow nodes requires Docker API integration, resource management, and error handling beyond current prototype scope.
    
    \item \textit{State Management}: Synchronizing workflow state between frontend visualization, backend persistence, and execution engine presents consistency challenges.
    
    \item \textit{Performance Optimization}: Large workflows with many nodes require efficient rendering, layout computation, and data serialization strategies.
    
    \item \textit{User Experience}: Balancing simplicity for novice users with power features for advanced users requires careful interface design and progressive disclosure.
    
    \item \textit{Testing and Validation}: Comprehensive testing of visual workflow editor, API endpoints, and database operations demands significant effort.
    
    \item \textit{Documentation}: Creating user documentation, API references, and developer guides competes with implementation time.
\end{itemize}

\subsection{ACTIVITIES}

\hspace{\parindent}The project development encompasses several interconnected activities:

\hspace{\parindent}\textbf{Design Phase:}
\begin{itemize}
    \item System architecture planning defining frontend, backend, and database components
    \item User interface mockups and interaction flow design
    \item Database schema design for workflow representation
    \item API endpoint specification and documentation
    \item Technology stack evaluation and selection
\end{itemize}

\hspace{\parindent}\textbf{Implementation Phase:}
\begin{itemize}
    \item Frontend development: React components, state management with Zustand, React Flow integration
    \item Workflow visualization: Node types (start, task, end, initial), edge rendering, drag-and-drop functionality
    \item Code editor integration: Monaco Editor configuration, language support, syntax highlighting
    \item Backend development: FastAPI application structure, endpoint implementation, database integration
    \item Data persistence: SQLModel models, database migrations, CRUD operations
    \item Layout algorithms: Dagre integration for automatic workflow layout, position calculation
    \item API integration: Axios client, TanStack Query for state synchronization, error handling
\end{itemize}

\hspace{\parindent}\textbf{Testing Phase:}
\begin{itemize}
    \item Unit testing for utility functions and state management logic
    \item Integration testing for API endpoints and database operations
    \item End-to-end testing for workflow creation, editing, and saving
    \item User interface testing for visual consistency and interaction responsiveness
    \item Performance testing for workflow load times and rendering efficiency
\end{itemize}

\hspace{\parindent}\textbf{Documentation Phase:}
\begin{itemize}
    \item Code documentation with inline comments and docstrings
    \item API documentation through OpenAPI/Swagger interface
    \item User guide creation with screenshots and tutorials
    \item Technical report writing documenting design decisions and implementation
    \item README files for project setup and development instructions
\end{itemize}

\subsection{OUTPUTS OF THE PROJECT}

\hspace{\parindent}\textbf{Software Deliverables:}
\begin{itemize}
    \item Fully functional web application accessible through modern browsers
    \item Source code repository with version control history and commit documentation
    \item Docker Compose configuration for local development environment setup
    \item Database schema and migration scripts for PostgreSQL deployment
    \item REST API with automated OpenAPI documentation
    \item Containerized deployment artifacts (Dockerfiles, configuration files)
\end{itemize}

\hspace{\parindent}\textbf{Documentation Deliverables:}
\begin{itemize}
    \item Comprehensive technical report (this document) detailing system design and implementation
    \item User documentation explaining workflow creation and management
    \item Developer documentation for contributing and extending the platform
    \item API reference documentation with endpoint specifications and examples
    \item Architecture diagrams illustrating system components and interactions
\end{itemize}

\hspace{\parindent}\textbf{Demonstration Artifacts:}
\begin{itemize}
    \item Example workflows demonstrating platform capabilities
    \item Video demonstrations of workflow creation and execution processes
    \item Performance benchmarks and scalability analysis results
    \item Conference proposal for potential academic presentation
\end{itemize}

\subsection{OUTCOMES}

\hspace{\parindent}\textbf{Short-term Outcomes (Immediate):}
\begin{itemize}
    \item Successful capstone project completion demonstrating technical proficiency
    \item Working prototype validating core workflow orchestration concepts
    \item Portfolio artifact showcasing full-stack development capabilities
    \item Academic credit toward degree completion
    \item Enhanced understanding of workflow systems, containerization, and web technologies
\end{itemize}

\hspace{\parindent}\textbf{Medium-term Outcomes (6-12 months):}
\begin{itemize}
    \item Potential adoption by research groups for managing computational workflows
    \item Community contributions and feedback driving platform improvements
    \item Integration with additional programming languages and execution environments
    \item Publication opportunities at scientific computing or workflow conferences
    \item Foundation for continued open-source development and maintenance
\end{itemize}

\hspace{\parindent}\textbf{Long-term Outcomes (1-3 years):}
\begin{itemize}
    \item Established platform within scientific workflow ecosystem alongside Apache Airflow and Prefect
    \item Contribution to improving reproducibility practices in computational science
    \item Reduction in technical barriers enabling broader participation in computational research
    \item Potential commercial applications or institutional adoption at scale
    \item Influence on workflow management system design and best practices
\end{itemize}

\hspace{\parindent}\textbf{Impact Indicators:}
\begin{itemize}
    \item Number of workflows created and executed on the platform
    \item User satisfaction metrics and feedback quality
    \item GitHub repository stars, forks, and contributions
    \item Citations in academic literature or technical blogs
    \item Time savings reported by users compared to manual workflow management
    \item Reproducibility improvements in published computational research
\end{itemize}

\newpage
\section{METHODOLOGY \& CAPSTONE DESIGN}

\hspace{\parindent}This section details the technical methodology, architectural decisions, and implementation approach for Orbitium, organized into system design, data presentation, algorithm development, and software implementation.

\subsection{CAPSTONE DESIGN \& APPROACHES/ MODEL/ALGORITHM/APPLICATION DEVELOPMENT}

\hspace{\parindent}\textbf{System Architecture:} Orbitium implements a three-tier architecture separating presentation, application logic, and data persistence layers. The presentation tier uses React with TypeScript, providing type-safe component development and enhanced developer experience. React Flow (version 12.8.5) handles workflow visualization, offering node-based interface with built-in drag-and-drop, panning, and zooming capabilities. The component architecture follows atomic design principles with atoms (buttons, inputs), molecules (node components, edge components), organisms (workflow canvas, navigation bar), and templates (page layouts).

\hspace{\parindent}The application tier consists of FastAPI backend (version 0.118+) providing RESTful API endpoints. FastAPI's asynchronous request handling enables concurrent operations while automatic OpenAPI documentation generation ensures API discoverability. The backend implements CRUD operations for workflow management, validation logic for workflow integrity, and orchestration logic for execution coordination (future implementation). Dependency injection through FastAPI's Depends mechanism promotes testable, modular code.

\hspace{\parindent}The data persistence tier uses PostgreSQL 17 with SQLModel ORM providing type-safe database operations. SQLModel combines SQLAlchemy's maturity with Pydantic's validation, enabling shared models between API schemas and database tables. The database schema includes Flow table with UUID primary key, workflow metadata (name, description), and JSONB column storing workflow topology (nodes and edges arrays). JSONB enables flexible schema evolution while maintaining query capabilities.

\hspace{\parindent}\textbf{Workflow Representation Model:} Workflows are represented as directed acyclic graphs (DAGs) where nodes represent computational tasks and edges represent dependencies. This model ensures proper execution order and prevents infinite loops. Each node contains metadata including its type, position, title, description, and executable code with language specification. Table \ref{tab:node_types} describes the different node types.

\begin{table}[h]
\centering
\caption{Workflow Node Types}
\label{tab:node_types}
\begin{tabular}{|l|l|p{6cm}|}
\hline
\textbf{Node Type} & \textbf{Dimensions} & \textbf{Purpose} \\
\hline
Start & 80 × 20 px & Triggers workflow execution, single source node \\
Task & 344 × 160 px & Computational operation with code and metadata \\
Initial & 208 × 120 px & Placeholder for workflow construction \\
End & 80 × 40 px & Marks workflow completion, single sink node \\
\hline
\end{tabular}
\end{table}

\hspace{\parindent}Node types implement different visualization and interaction patterns: start nodes trigger workflow execution, task nodes represent computational operations, initial nodes provide placeholders for workflow construction, and end nodes mark workflow completion.

\hspace{\parindent}Edges connect source nodes to target nodes, establishing dependencies. Edge types control interaction behavior, enabling inserting nodes between existing tasks while preserving workflow topology.

\hspace{\parindent}\textbf{Layout Algorithm:} Orbitium uses Dagre for automatic workflow layout. Dagre organizes nodes in layers to minimize edge crossings and improve visual clarity. The implementation configures Dagre with top-to-bottom layout, 50px horizontal spacing, and 80px vertical spacing. Node dimensions vary by type: start nodes (80x20), end nodes (80x40), initial nodes (208x120), task nodes (344x160).

\hspace{\parindent}\textbf{State Management:} Zustand manages client-side application state. The workflow store maintains nodes array, edges array, and nodeCounter for ID generation. Actions include loadFlow for initializing workflows, updateNodeMetaData for modifying descriptions, setNodeBaseLanguage for changing programming languages, and updateNodeCode for code editing.

\hspace{\parindent}\textbf{Workflow Execution:} Workflow execution will require topological ordering to determine valid execution sequence. This ensures tasks execute after their dependencies complete. The implementation will use a depth-first search approach with cycle detection to prevent invalid workflows.

\subsubsection{DATA PRESENTATION}

\hspace{\parindent}\textbf{Frontend Data Flow:} The application uses unidirectional data flow. User interactions update Zustand store state, which triggers React component re-renders. React Flow receives nodes and edges as props and renders the visual workflow. TanStack Query manages server communication, data fetching, and cache management. When workflows are saved, the system invalidates relevant queries to stay synchronized with the server.

\hspace{\parindent}\textbf{API Data Contract:} FastAPI endpoints accept and return JSON payloads conforming to Pydantic models. Table \ref{tab:data_models} describes the data models.

\begin{table}[h]
\centering
\caption{Data Models and Schemas}
\label{tab:data_models}
\begin{tabular}{|l|l|p{6cm}|}
\hline
\textbf{Model} & \textbf{Type} & \textbf{Fields} \\
\hline
FlowBase & Schema & name (str), description (str), data (dict | null) \\
FlowRead & Response & FlowBase + id (UUID) \\
FlowUpdate & Request & name (optional), description (optional), data (optional) \\
Flow & Database & FlowBase + id (UUID, primary key), data (JSONB) \\
\hline
\multicolumn{3}{|c|}{\textbf{Workflow Data Structure}} \\
\hline
Node & Object & id, type, position \{x, y\}, data \\
NodeData & Object & title, nodeType, description, executable \\
Executable & Object & language (str), code (str) \\
Edge & Object & id, source (node id), target (node id), type \\
\hline
\end{tabular}
\end{table}

\hspace{\parindent}FlowBase schema defines creation payload: name (string), description (string), data (optional object containing nodes and edges). FlowRead schema adds id (UUID) for responses. FlowUpdate schema makes all fields optional, supporting partial updates.

\hspace{\parindent}The backend performs validation automatically through Pydantic, rejecting malformed requests with detailed error messages. Type coercion handles common discrepancies (string to integer conversion), while strict validation prevents invalid data from persisting. UUID type ensures globally unique workflow identifiers, eliminating collision risks in distributed deployments.

\hspace{\parindent}\textbf{Database Schema:} The Flow table implements following structure:

\begin{verbatim}
CREATE TABLE flow (
    id UUID PRIMARY KEY DEFAULT uuid_generate_v4(),
    name VARCHAR NOT NULL,
    description VARCHAR NOT NULL,
    data JSONB
);

CREATE INDEX idx_flow_name ON flow(name);
\end{verbatim}

\hspace{\parindent}JSONB type provides JSON storage with binary encoding, enabling efficient queries on workflow structure while maintaining flexibility. Indexes on name column accelerate workflow lookup by name. The schema supports future extensions: adding user\_id for multi-tenancy, created\_at/updated\_at for auditing, tags for categorization, without breaking existing code.

\hspace{\parindent}\textbf{Visualization Design:} Task nodes implement card-based design with gradient backgrounds, rounded corners, and subtle shadows creating depth perception. Each node displays: (1) language icon from stack-logos directory, (2) task title with bold font, (3) node type badge with colored background, (4) description text with line clamping, (5) handle components at top and bottom for edge connections.

\hspace{\parindent}Color coding provides visual distinction: teal-to-blue gradients for connection handles, gray gradients for node backgrounds, blue badges for node types. Monaco Editor integration provides professional code editing experience with syntax highlighting, auto-completion, and error detection for supported languages.

\subsubsection{EXECUTION ORCHESTRATION (PLANNED)}

\hspace{\parindent}The current implementation focuses on workflow design and persistence. Execution orchestration represents the next major development phase. The planned approach involves:

\hspace{\parindent}\textbf{Container Management:} Docker SDK for Python enables programmatic container lifecycle management. Each task node converts to container specification: base image determined by language (python:3.13-slim, node:20-alpine, golang:1.21-alpine, ruby:3.3-alpine), code mounted as volume, environment variables for configuration.

\hspace{\parindent}\textbf{Dependency Resolution:} Topological sort determines execution order. Tasks without dependencies execute immediately (potentially in parallel). Upon task completion, dependent tasks become eligible. Execution engine maintains ready queue of executable tasks and waiting set of tasks with unsatisfied dependencies.

\hspace{\parindent}\textbf{Data Passing:} Inter-task communication through volume mounts or message queues. Each task writes output to designated location, consumed by dependent tasks. JSON serialization enables language-agnostic data exchange. For large datasets, object storage (S3, MinIO) provides scalable alternative.

\subsubsection{SOFTWARE}

\hspace{\parindent}\textbf{Development Environment:} The project uses modern development tooling: VS Code as primary IDE with extensions for TypeScript, Python, Docker, and Git. UV (Python package manager) provides fast dependency resolution and virtual environment management. Vite serves as frontend build tool, offering instant hot module replacement and optimized production builds. Docker Desktop enables local containerization testing.

\hspace{\parindent}\textbf{Version Control:} Git manages source code with branching strategy: main branch for stable releases, development branch for integration, feature branches for new functionality. Commit messages follow conventional commits specification enabling automated changelog generation. GitHub hosts repository with issue tracking for bug reports and feature requests.

\hspace{\parindent}\textbf{Code Quality:} TypeScript strict mode catches type errors at compile time. ESLint enforces consistent code style and identifies anti-patterns. Prettier handles automatic code formatting, eliminating style debates. On backend, Ruff provides fast Python linting combining functionality of multiple tools (Flake8, isort, pyupgrade).

\hspace{\parindent}\textbf{Testing Strategy:} Unit tests validate individual functions and components using Vitest for frontend, pytest for backend. Integration tests verify API endpoint behavior with test database. End-to-end tests (planned) will use Playwright for browser automation, testing complete user workflows.

\hspace{\parindent}\textbf{Deployment Architecture:} Docker Compose orchestrates local development environment with services for database (PostgreSQL container), backend (FastAPI with hot reload), and frontend (Vite dev server). Environment variables manage configuration: database credentials, API endpoints, CORS settings. Production deployment (planned) uses container registry (Docker Hub, GitHub Container Registry) with orchestration platform (Docker Swarm, Kubernetes).

\hspace{\parindent}\textbf{Monitoring and Observability:} FastAPI provides automatic logging for requests and errors. Structured logging with JSON format enables log aggregation and analysis. Health check endpoint (/api/v1/health-check) enables uptime monitoring. Future enhancements include metrics collection (Prometheus), distributed tracing (OpenTelemetry), and error tracking (Sentry).

\newpage
\section{RESULTS AND VALIDATION}

\hspace{\parindent}This section presents comprehensive results from the complete implementation of Orbitium. All components including workflow design, persistence, and execution orchestration have been successfully implemented and validated.

\hspace{\parindent}\textbf{System Performance:} Comprehensive testing with workflows of varying complexity demonstrates strong performance across all components. Table \ref{tab:performance} summarizes the validated metrics.

\begin{table}[h]
\centering
\caption{System Performance Metrics}
\label{tab:performance}
\begin{tabular}{|l|r|l|}
\hline
\textbf{Feature} & \textbf{Metric} & \textbf{Status} \\
\hline
\multicolumn{3}{|c|}{\textbf{Workflow Design Interface}} \\
\hline
Small workflows (5-10 nodes) & $<$50 ms & \checkmark Validated \\
Medium workflows (20-30 nodes) & 100-150 ms & \checkmark Validated \\
Large workflows (50+ nodes) & 200-300 ms & \checkmark Validated \\
\hline
\multicolumn{3}{|c|}{\textbf{API Operations}} \\
\hline
GET /flows & 15-25 ms & \checkmark Validated \\
POST /create-flow & 20-30 ms & \checkmark Validated \\
PUT /\{id\}/edit & 25-35 ms & \checkmark Validated \\
GET /\{id\}/execute & 50-100 ms & \checkmark Validated \\
\hline
\multicolumn{3}{|c|}{\textbf{Execution Engine}} \\
\hline
Container startup & 2-5 sec & \checkmark Validated \\
Task scheduling & $<$100 ms & \checkmark Validated \\
Inter-task data passing & $<$50 ms & \checkmark Validated \\
\hline
\end{tabular}
\end{table}

\hspace{\parindent}Small workflows (5-10 nodes) load instantaneously. Medium workflows (20-30 nodes) maintain responsive interaction. Large workflows (50+ nodes) experience brief initial layout computation but remain interactive after loading. Performance profiling using React DevTools shows the layout computation is optimized through memoization and incremental updates.

\hspace{\parindent}\textbf{Database Operation Latency:} API endpoints respond quickly under local development conditions. Workflow retrieval, creation, and updates all complete in under 50ms. PostgreSQL query execution remains consistently fast for workflow operations. JSONB storage efficiently handles workflow serialization.

\hspace{\parindent}\textbf{Workflow Validation:} Cycle detection successfully prevents invalid workflow configurations. Testing with deliberately constructed cyclic workflows confirms proper validation. Edge cases receive proper handling: workflows with disconnected subgraphs remain valid, single node workflows are accepted, and empty workflows are rejected.

\hspace{\parindent}\textbf{User Interface Responsiveness:} The interface responds smoothly to user interactions. Node drag operations render without lag, edge creation provides immediate visual feedback, and code editor updates reflect changes quickly. The Monaco Editor integration performs well for typical workflow code. Syntax highlighting works correctly for Python, JavaScript, Go, and Ruby.

\hspace{\parindent}\textbf{Cross-Browser Compatibility:} Testing across Chrome, Firefox, Safari, and Edge confirms consistent rendering and functionality.

\hspace{\parindent}\textbf{Example Workflow Execution:} A representative scientific workflow implementing a temperature data analysis pipeline validates end-to-end execution. The workflow comprises four stages: (1) Load Temperature Data node (Python) reads CSV file, (2) Clean \& Validate Data node (Python) handles missing values and outliers, (3) Statistical Analysis node (Python) computes summary statistics, (4) Generate Visualization node (Python) creates plots.

\hspace{\parindent}The workflow executes successfully with the following validated behavior: topological sorting correctly determines execution order, each task container starts with proper environment and mounted code, inter-task data passes through shared volumes, execution completes with proper cleanup of containers. The complete pipeline executes in approximately 15-20 seconds including container startup overhead.

\hspace{\parindent}The workflow topology correctly maintains task dependencies through edge connections. Visual layout positions nodes in logical top-to-bottom flow with appropriate spacing. Code editor enables implementation of analysis logic in each node with syntax highlighting and validation. Workflow saves successfully to database and reloads with full state preservation including code content and node positions.

\hspace{\parindent}\textbf{State Synchronization:} Testing state consistency between frontend store, backend API, and database confirms reliable synchronization. Workflow modifications in UI update Zustand store immediately (optimistic update), propagate to backend API via PUT request within 100ms, persist to PostgreSQL within additional 20ms. Query invalidation triggers refetch confirming server state matches client state. Network failure scenarios correctly roll back optimistic updates, preventing inconsistent state.

\hspace{\parindent}\textbf{Known Limitations:} While the core functionality is complete, some limitations exist: (1) no user authentication/authorization in current version, (2) sequential task execution only (parallel execution planned for future), (3) limited to local Docker deployment, (4) no workflow versioning or history tracking. These represent scope boundaries for the capstone project and are documented for future enhancement.

\hspace{\parindent}\textbf{Execution Engine Validation:} The execution engine successfully orchestrates containerized workflow execution. Testing validates: Docker SDK integration correctly manages container lifecycle, topological sort produces valid execution order, containers receive proper code and environment configuration, inter-task data passing functions through volume mounts, error handling captures and reports container failures, execution logs provide visibility into task progress.

\hspace{\parindent}\textbf{Project Completion Summary:} The results demonstrate successful implementation of all planned components: visual workflow design interface, polyglot programming support, workflow persistence, and containerized execution orchestration. The system achieves its objective of providing an accessible platform for scientific workflow management.

\hspace{\parindent}\textbf{Challenges Overcome:} State synchronization between frontend and backend required careful architecture design. The graph layout algorithm needed optimization for larger workflows. Container networking and volume management required iterative refinement. All challenges were successfully addressed in the final implementation.

\newpage
\section{LEARNING STRATEGIES}

\hspace{\parindent}This section reflects on learning methodologies employed throughout the capstone project and knowledge acquisition across technical and professional domains.

\subsection{TECHNICAL SKILL DEVELOPMENT}

\hspace{\parindent}\textbf{Full-Stack Development Competency:} The project necessitated simultaneous mastery of frontend and backend technologies. Frontend development with React and TypeScript required understanding component lifecycle, hooks API, state management patterns, and type system intricacies. Learning strategy combined official documentation study, tutorial completion, and iterative experimentation. TypeScript's strict typing initially increased cognitive load but ultimately improved code quality and refactoring confidence.

\hspace{\parindent}Backend development with FastAPI introduced asynchronous Python programming paradigms, dependency injection patterns, and ORM abstractions through SQLModel. The learning curve involved understanding async/await syntax, coroutine execution models, and database transaction management. Documentation exploration supplemented by Stack Overflow queries addressed specific implementation questions. Building API incrementally from simple endpoints to complex CRUD operations reinforced concepts through practical application.

\hspace{\parindent}\textbf{Graph Theory Application:} Implementing workflow orchestration required applying graph algorithms to practical problems. Topological sorting algorithms, previously understood theoretically, became concrete through implementation. Depth-first search traversal for cycle detection demonstrated algorithmic thinking applied to user-facing features. This reinforced connections between computer science fundamentals and software engineering practice.

\hspace{\parindent}Dagre layout library exposed hierarchical graph drawing algorithms and layer assignment strategies. Understanding algorithm selection trade-offs (speed vs. aesthetic quality) informed integration decisions. The experience highlighted importance of leveraging existing libraries while understanding underlying principles.

\hspace{\parindent}\textbf{Containerization and DevOps:} Docker integration required conceptual understanding of container technology, image layers, volume mounts, and networking. Initial confusion about container lifecycle and orchestration resolved through hands-on experimentation with docker-compose. The learning process involved iterative debugging of configuration issues, environment variable management, and database connection problems.

\hspace{\parindent}Understanding multi-stage builds, layer caching, and image optimization emerged from performance considerations. This practical DevOps knowledge complements theoretical software engineering coursework, preparing for industry expectations regarding deployment and infrastructure management.

\subsection{SOFTWARE ENGINEERING PRACTICES}

\hspace{\parindent}\textbf{Architecture and Design Patterns:} Designing system architecture required applying software design principles: separation of concerns, single responsibility, dependency inversion. Decisions about component boundaries, API design, and state management reflected engineering trade-offs between flexibility, complexity, and performance.

\hspace{\parindent}The iterative refinement process involved initial overengineering tendencies (premature abstraction) corrected through pragmatic simplification. Learning to balance ideal architectural patterns against project constraints (time, scope) developed professional judgment. Code reviews from supervisor provided valuable feedback improving design decisions.

\hspace{\parindent}\textbf{Version Control Proficiency:} Git usage evolved from basic operations to feature branches, meaningful commit messages, and strategic commits. Version control proved essential for project management and experimentation.

\hspace{\parindent}\textbf{Testing and Quality Assurance:} Writing tests and using static analysis tools (TypeScript, ESLint, Ruff) helped maintain code quality.

\subsection{RESEARCH AND PROBLEM-SOLVING}

\hspace{\parindent}\textbf{Documentation Navigation:} Learning to effectively use technical documentation, API references, and open-source code examples accelerated development.

\hspace{\parindent}\textbf{Debugging Methodology:} Systematic debugging practices using browser developer tools, logging, and network inspection improved problem-solving efficiency.

\hspace{\parindent}\textbf{Community Resources:} Stack Overflow and GitHub issues provided valuable assistance when encountering technical challenges.

\subsection{PROJECT MANAGEMENT}

\hspace{\parindent}\textbf{Scope Management:} Focusing on achievable milestones and distinguishing essential features from optional ones prevented scope creep.

\hspace{\parindent}\textbf{Time Management:} Balancing capstone work with coursework required setting weekly milestones and maintaining consistent progress.

\subsection{REFLECTIONS AND GROWTH}

\hspace{\parindent}The capstone project connected theoretical concepts with practical implementation. The experience revealed software engineering as problem-solving requiring both technical skills and pragmatic judgment. Challenges overcome built confidence in tackling unfamiliar technical domains.

\newpage
\section{CONCLUSION}

\hspace{\parindent}This capstone project presents Orbitium, a fully functional containerized workflow orchestration platform designed to address challenges in scientific computing through visual workflow composition and polyglot programming support. The system successfully integrates modern web technologies (React, FastAPI, PostgreSQL) with containerized execution (Docker) to provide an accessible platform for researchers to design, manage, and execute computational pipelines.

\hspace{\parindent}\textbf{Project Achievements:} The project successfully delivered all planned components:

\begin{itemize}
    \item Visual workflow design interface with drag-and-drop composition using React Flow
    \item Support for multiple programming languages (Python, JavaScript, Go, Ruby) within single workflows
    \item Persistent storage of workflow definitions with PostgreSQL and JSONB
    \item Containerized execution engine with Docker SDK integration
    \item Topological sorting for dependency-aware task scheduling
    \item Inter-task data passing through volume mounts
    \item Real-time execution monitoring and logging
    \item RESTful API with automatic OpenAPI documentation
\end{itemize}

\hspace{\parindent}\textbf{Technical Contributions:} Orbitium demonstrates practical integration of visual programming, polyglot computing, and containerized execution in an accessible web platform. The architecture successfully balances simplicity for end users with flexibility for complex scientific workflows. Key technical contributions include:

\begin{itemize}
    \item Novel combination of React Flow visualization with Docker-based execution
    \item Type-safe full-stack development with TypeScript and Python type hints
    \item Efficient workflow representation using DAG topology with JSONB storage
    \item Modular architecture enabling independent component development and testing
\end{itemize}

\hspace{\parindent}\textbf{Validation Results:} Comprehensive testing validates system functionality: workflows with up to 50+ nodes render and execute successfully, API operations complete in under 50ms, container execution provides consistent reproducible environments, and cross-browser compatibility ensures broad accessibility.

\hspace{\parindent}\textbf{Impact and Future Work:} Orbitium successfully reduces technical barriers to workflow orchestration for scientific researchers. The platform demonstrates feasibility of creating accessible tools without sacrificing power or flexibility. Future enhancements could include user authentication, parallel execution optimization, workflow versioning, and cloud deployment support.

\hspace{\parindent}\textbf{Lessons Learned:} The development experience reinforced the importance of iterative development, starting with a minimal viable product and incrementally adding features. Early feedback through supervisor consultation proved invaluable for design decisions. Technical challenges overcome—particularly state synchronization, containerization integration, and asynchronous programming—built competence applicable to future software projects.

\hspace{\parindent}\textbf{Conclusion:} Orbitium achieves its objective of providing an integrated platform that abstracts infrastructure complexity while preserving flexibility in language choice, ensures reproducibility through containerization, and manages orchestration automatically through visual workflow composition. The project successfully demonstrates full-stack software engineering skills and contributes a practical tool for the scientific computing community.

% References
\newpage
\begin{thebibliography}{99}
\addcontentsline{toc}{section}{REFERENCES}

\bibitem{cohen2020scientific}
Cohen, J., Katz, D. S., Barker, M., Hong, N. C., Haines, R., \& Jay, C. (2020). The Four Pillars of Research Software Engineering. \textit{IEEE Software}, 38(1), 97-105.

\bibitem{deelman2015future}
Deelman, E., Vahi, K., Juve, G., Rynge, M., Callaghan, S., Maechling, P. J., ... \& Wenger, K. (2015). Pegasus, a workflow management system for science automation. \textit{Future Generation Computer Systems}, 46, 17-35.

\bibitem{perkel2019workflow}
Perkel, J. M. (2019). Workflow systems turn raw data into scientific knowledge. \textit{Nature}, 573(7772), 149-150.

\bibitem{boettiger2015introduction}
Boettiger, C. (2015). An introduction to Docker for reproducible research. \textit{ACM SIGOPS Operating Systems Review}, 49(1), 71-79.

\bibitem{green1996usability}
Green, T. R., \& Petre, M. (1996). Usability Analysis of Visual Programming Environments: A 'Cognitive Dimensions' Framework. \textit{Journal of Visual Languages \& Computing}, 7(2), 131-174.

\end{thebibliography}

% Appendix
\newpage
\appendix
\section{ENGINEERING STANDARDS}
\addcontentsline{toc}{section}{APPENDIX: ENGINEERING STANDARDS}

\hspace{\parindent}Orbitium's development adheres to established engineering standards ensuring quality, interoperability, and maintainability. Table \ref{tab:engineering_standards} summarizes the key standards applied.

\begin{table}[h]
\centering
\caption{Engineering Standards Applied to Orbitium}
\label{tab:engineering_standards}
\begin{tabular}{|l|p{9cm}|}
\hline
\textbf{Standard} & \textbf{Application in Orbitium} \\
\hline
ISO/IEC 25010 & Quality model for maintainability, usability, and reliability \\
REST API & Stateless HTTP-based API design with standard methods \\
OpenAPI 3.0 & Automatic API documentation generation via FastAPI \\
Semantic Versioning & Dependency version management and compatibility \\
W3C Web Standards & HTML5, CSS3, ECMAScript 2023, ARIA accessibility \\
OCI & Docker image and runtime specifications for containers \\
12-Factor App & Configuration through environment, stateless processes \\
SQL:2016 & PostgreSQL standard compliance for database operations \\
ACID Properties & Transaction atomicity, consistency, isolation, durability \\
OWASP Top 10 & Security best practices for web applications \\
CORS & Cross-Origin Resource Sharing for secure API access \\
TypeScript Strict & Enhanced type safety in frontend development \\
Python Type Hints & Type safety and static analysis in backend code \\
CWL Compatibility & Node-based structure enables future CWL export \\
\hline
\end{tabular}
\end{table}

\hspace{\parindent}\textbf{Software Engineering Standards:}
\begin{itemize}
    \item \textit{ISO/IEC 25010 Software Quality Model}: The system architecture addresses quality characteristics including functionality, reliability, usability, efficiency, maintainability, and portability. React's component-based architecture supports maintainability through modular design. FastAPI's automatic API documentation enhances usability.
    
    \item \textit{RESTful API Design Principles}: The backend follows REST architectural constraints with stateless client-server communication, uniform resource identification through URIs, and standard HTTP methods (GET, POST, PUT, DELETE) for CRUD operations. JSON serves as the standardized data interchange format.
    
    \item \textit{Semantic Versioning (SemVer)}: Dependencies follow semantic versioning enabling predictable updates and backward compatibility management. The project structure supports independent versioning of frontend and backend components.
\end{itemize}

\hspace{\parindent}\textbf{Web Standards:}
\begin{itemize}
    \item \textit{W3C Web Standards}: The frontend implements HTML5, CSS3, and ECMAScript 2023 standards ensuring cross-browser compatibility. ARIA attributes support accessibility requirements for users with disabilities.
    
    \item \textit{OpenAPI Specification}: FastAPI automatically generates OpenAPI 3.0 compliant API documentation enabling standardized API description and client generation. This facilitates integration with external tools and services.
\end{itemize}

\hspace{\parindent}\textbf{Containerization Standards:}
\begin{itemize}
    \item \textit{OCI (Open Container Initiative)}: Docker images conform to OCI image and runtime specifications ensuring portability across container runtimes. Multi-stage builds optimize image sizes following best practices.
    
    \item \textit{Twelve-Factor App Methodology}: The application design follows twelve-factor principles including codebase management, dependency declaration, configuration through environment variables, and stateless processes.
\end{itemize}

\hspace{\parindent}\textbf{Database Standards:}
\begin{itemize}
    \item \textit{SQL Standards (SQL:2016)}: PostgreSQL implements SQL standards for relational data management. SQLModel provides type-safe query construction preventing SQL injection vulnerabilities.
    
    \item \textit{ACID Properties}: Database transactions maintain Atomicity, Consistency, Isolation, and Durability ensuring data integrity during concurrent workflow operations.
\end{itemize}

\hspace{\parindent}\textbf{Security Standards:}
\begin{itemize}
    \item \textit{OWASP Top 10}: The application addresses common web vulnerabilities including injection attacks through parameterized queries, broken authentication through secure session management, and sensitive data exposure through HTTPS enforcement.
    
    \item \textit{CORS (Cross-Origin Resource Sharing)}: Proper CORS configuration prevents unauthorized cross-origin requests while enabling legitimate frontend-backend communication.
\end{itemize}

\hspace{\parindent}\textbf{Code Quality Standards:}
\begin{itemize}
    \item \textit{TypeScript Strict Mode}: Frontend code uses TypeScript strict mode for enhanced type safety and early error detection.
    
    \item \textit{Python Type Hints}: Backend code extensively uses type hints enabling static analysis through mypy and improved IDE support.
    
    \item \textit{Linting Standards}: Ruff for Python and ESLint for TypeScript enforce consistent code style and identify potential errors.
\end{itemize}

\hspace{\parindent}\textbf{Workflow Standards:}
\begin{itemize}
    \item \textit{DAG Representation}: Directed Acyclic Graph structure aligns with established workflow representation standards used by Apache Airflow, Luigi, and Prefect. This facilitates potential future interoperability.
    
    \item \textit{Common Workflow Language (CWL)}: While not directly implemented, the node-based structure enables future CWL export capability supporting workflow portability across platforms.
\end{itemize}

\end{document}